\documentclass[12pt,oneside]{article}

\usepackage[margin=1in]{geometry}
\usepackage{times}
\usepackage[hidelinks]{hyperref}

\pagestyle{empty}

\begin{document}

\begin{flushleft}

MISJustice Alliance \\
On behalf of Elvis Ryland Nuno and other victims \\
\today

\vspace{0.3in}

Federal Bureau of Investigation \\
Civil Rights Division, Seattle Field Office \\
1110 3rd Avenue \\
Seattle, WA 98101

\vspace{0.3in}

\textbf{RE: Comprehensive Criminal Investigation Referral} \\
\textbf{Pattern of Institutional Corruption, RICO Violations, Civil Rights Conspiracy} \\
\textbf{Danielle Chard, E'Lise Chard, YWCA Missoula, Missoula PD, Associated Officials} \\
\textbf{Federal Jurisdiction: 18 U.S.C. § 1962 (RICO); 42 U.S.C. § 1983, 1985 (Civil Rights)} \\
\textbf{Interstate Criminal Activity: Washington State \& Montana (2015--2025)}

\vspace{0.3in}

Dear Special Agent in Charge:

\vspace{0.2in}

MISJustice Alliance respectfully submits this \textbf{comprehensive criminal investigation referral} documenting a decade-long pattern of institutional corruption, RICO violations, and systematic civil rights abuses spanning Washington State and Montana. This submission requests immediate investigation by the Seattle Field Office Civil Rights squad, in coordination with the U.S. Attorney's Offices for the District of Montana and the Western District of Washington and the Department of Justice Criminal Division's Public Integrity Section.

\section*{SCOPE OF INVESTIGATION REQUEST}

While the majority of submitted documentation focuses on the systematic victimization of \textbf{Elvis Ryland Nuno} over 11+ years (2015--present), the evidence establishes \textbf{beyond reasonable doubt} that individual criminal conduct represents merely the visible manifestation of endemic institutional corruption affecting numerous victims. We request investigation targeting:

\begin{enumerate}
\item \textbf{Individual Perpetrators:} Danielle Christine Chard, E'Lise Marie Chard
\item \textbf{Institutional Enterprise:} YWCA Missoula and governance structure
\item \textbf{Law Enforcement Actors:} Detective Connie Brueckner, Officer Ethan Smith, MPD supervisors
\item \textbf{Prosecutorial Actors:} Missoula County and King County prosecutors
\item \textbf{Pattern Investigation:} Expanded inquiry into institutional corruption affecting multiple victims (2012--2025)
\end{enumerate}

\section*{FEDERAL JURISDICTION}

\subsection*{RICO Violations (18 U.S.C. § 1962)}

\textbf{Enterprise:} YWCA Missoula (nonprofit corporation), Missoula Police Department, Missoula County Prosecutor's Office, coordinated officials across Washington and Montana.

\textbf{Pattern of Racketeering:} 20+ predicate acts documented across 2012--2025:

\begin{itemize}
\item Mail/wire fraud: Federal grant applications, sponsor reporting, interstate communications.
\item Obstruction of justice: Evidence suppression, complaint suppression, cover-ups.
\item HIPAA violations: Criminal disclosure of protected health information (42 U.S.C. § 1320d-6).
\item Witness tampering: Intimidation of complainants and potential witnesses.
\item Interstate criminal activity: Coordinated false prosecution across WA/MT borders.
\end{itemize}

\textbf{Interstate Commerce Impact:} Federal funding fraud, multi-state conspiracy, wire communications, denial of federally-protected civil rights.

\textbf{Damages:} \$60+ million across documented victims:

\begin{itemize}
\item Nuno case: \$18M+ (including RICO treble damages).
\item Institutional victims: \$42.8M+ (conservative estimate, multiple victims).
\end{itemize}

\subsection*{Civil Rights Violations (42 U.S.C. § 1983, 1985)}

\textbf{Pattern and Practice:} Systematic deprivation of constitutional rights under color of law:

\begin{itemize}
\item \textbf{First Amendment:} Retaliation for protected speech (YWCA complaint, public criticism).
\item \textbf{Fourth Amendment:} Warrantless home invasions, unreasonable searches/seizures.
\item \textbf{Eighth Amendment:} Excessive bail (\$50,000--\$100,000 for misdemeanors).
\item \textbf{Fourteenth Amendment:} Due process violations, malicious prosecution.
\end{itemize}

\textbf{Conspiracy:} Coordinated actions between YWCA, law enforcement, and prosecutors to deprive citizens of equal protection and constitutional rights.

\section*{EVIDENCE BEYOND NUNO CASE: ENDEMIC CORRUPTION}

The enclosed \textbf{YWCA Institutional Corruption Supplemental} (Document \#11) establishes RICO predicate acts \textbf{independent of the Nuno case}, demonstrating a broader criminal enterprise that has harmed multiple victims and families over more than a decade.

\subsection*{Multiple Documented Victims}

\begin{itemize}
\item \textbf{Arthur Brown:} YWCA Missoula submitted defamatory letters to a court recommending permanent restraining orders without adequate factual basis, resulting in separation of a father from his children through fraud upon the court.
\item \textbf{Anonymous Client \#1:} HIPAA violations in which counseling information was leaked to an abusive ex-partner through staff’s personal connections, coupled with documented ADA discrimination and failure to accommodate disabilities.
\item \textbf{Anonymous Client \#2:} Discriminatory service provision based on perceived client “compliance” rather than need, alongside fraudulent reporting to federal sponsors claiming successful outcomes that were not delivered.
\item \textbf{Multiple Employees (2023--2024 reviews):} Reports of “racist undertones,” untrained staff, exploitative working conditions, and organizational dysfunction creating dangerous conditions for vulnerable clients.
\end{itemize}

\subsection*{Institutional Governance Capture}

\textbf{Detective Connie Brueckner} simultaneously served as:

\begin{itemize}
\item A YWCA Missoula Board Member at-large.
\item The lead Missoula Police Department investigator of complaints against YWCA employees.
\item A primary law enforcement actor in constitutional violations against complainants, including Mr. Nuno.
\end{itemize}

This unprecedented conflict of interest enabled systematic abuse, undisclosed Brady material, and functional immunity from accountability for YWCA staff and associated officers, as complainants could not obtain an independent investigation of their grievances.

\subsection*{Federal Oversight Context}

\textbf{2012--2016 DOJ Investigation:} Missoula was identified as the “rape capital of the country” due to systemic failures in sexual-assault response. YWCA Missoula was designated as a key “community partner” in the resulting reforms, creating additional conflicts when YWCA board members, including a sitting detective, controlled the response to complaints against YWCA staff.

\textbf{2023 LifeGuard Group Matter:} A Montana DOJ Agent documented “concerning” strip search allegations involving a survivor and child at a facility run by the Sheriff’s Office Chaplain (receiving over \$500,000 from the Gianforte Foundation). Despite the severity of the allegations, no criminal investigation was launched and safe-house operations remained effectively unregulated.

\subsection*{Institutional Refusal to Respond to Formal Notices}

MISJustice Alliance sent formal legal notices to YWCA Missoula on \textbf{October 12, 2025} and \textbf{December 5, 2025}, documenting:

\begin{itemize}
\item Detective Brueckner’s dual law enforcement/board role and resulting conflict of interest.
\item Patterns of retaliation and confidential information misuse against complainants, including Mr. Nuno.
\item E'Lise Chard’s use of YWCA employment to coordinate and amplify a harassment campaign.
\item Recruitment of YWCA client Tyleen Root for defamation and harassment directed at Mr. Nuno.
\item Destruction of Mr. Nuno’s ability to safely shelter vulnerable women in his personal network.
\item Federal grant fraud indicators and mechanisms of institutional capture.
\end{itemize}

The December 5 notice was copied to senior YWCA staff, the Missoula Mayor, Missoula Police Chief, County Attorney, and other key officials. \textbf{To date, YWCA Missoula has provided no response to either notice.}\footnote{See Evidentiary Documentation \S{} 2.9--2.10, \texttt{Formal-Complaints/ywca-first-notice-10-12-2025.pdf} and \texttt{ywca-second-notice-12-5-2025.pdf}.} This continued institutional silence, following detailed legal allegations copied to governmental oversight, establishes:

\begin{itemize}
\item Consciousness of institutional guilt.
\item Deliberate indifference to civil rights violations (supporting \textit{Monell} liability).
\item A pattern of suppressing complaints (mirroring the non-response to Mr. Nuno’s 2018 formal complaint).
\item An institutional policy of non-accountability enabling ongoing constitutional violations.
\end{itemize}

The absence of \textit{any} response---not even an acknowledgment---to formal legal notices underscores the depth of institutional capture and the futility of internal remedies, necessitating federal intervention.

\section*{REQUESTED FEDERAL ACTIONS}

\begin{enumerate}
\item \textbf{Civil Rights Investigation:} Comprehensive FBI investigation by the Seattle Field Office Civil Rights squad into pattern and practice violations, institutional conspiracy, and interstate criminal coordination.
\item \textbf{Federal Grand Jury:} Coordination with the U.S. Attorney, District of Montana and U.S. Attorney, Western District of Washington to impanel grand juries for RICO and civil rights prosecutions.
\item \textbf{Referral to Public Integrity Section:} Close coordination with the DOJ Criminal Division Public Integrity Section regarding potential public corruption and prosecutorial misconduct by Missoula County and King County officials.
\item \textbf{HHS Office for Civil Rights:} HIPAA criminal investigation and penalties.
\item \textbf{Federal Grant Audit:} DOJ/HHS investigation of grant fraud and misrepresentation to federal agencies.
\item \textbf{Asset Forfeiture:} 18 U.S.C. § 1963 forfeiture of RICO enterprise proceeds.
\item \textbf{Witness Protection:} Consideration for victims facing retaliation threats.
\end{enumerate}

\section*{SUBMISSION CONTENTS (11 Documents)}

\begin{enumerate}
\item Case Summary.
\item This FBI Seattle Civil Rights Cover Letter.
\item FBI Montana Field Office Civil Rights Cover Letter (Montana-focused).
\item Montana Attorney General Cover Letter.
\item Western Washington Attorney General Cover Letter.
\item U.S. Attorney, District of Montana Cover Letter.
\item U.S. Attorney, Western District of Washington Cover Letter.
\item DOJ Public Integrity Section Cover Letter.
\item Danielle Chard Criminal Report.
\item E'Lise Chard Criminal Report.
\item YWCA Institutional Corruption Supplemental, Elvis Nuno Sworn Declaration, and comprehensive evidentiary documentation.
\end{enumerate}

\section*{EVIDENTIARY STANDARD AND URGENCY}

[Retain your existing evidentiary standard and urgency paragraphs here.]

We respectfully request an initial meeting or intake conference with the Seattle Field Office Civil Rights team to discuss investigation scope, coordination with the two U.S. Attorney's Offices, and referral to the Public Integrity Section. The submitted evidence is comprehensive and ready for immediate investigative triage.

\vspace{0.3in}

\noindent
Respectfully submitted,

\vspace{0.3in}

\noindent
MISJustice Alliance \\
On behalf of Elvis Ryland Nuno and other victims

\vspace{0.2in}

\noindent
Contact: [\textit{Contact information will be provided only through MISJustice Alliance after confirmation of concrete plan to prevent any retaliation or adverse action against reporting victims and witnesses.}]

\end{flushleft}

\end{document}
